%% ma: L12-B17-0024
%% lop: 12
%% chuong: 5
%% bai: 17
%% dang: 12.17.1
%% loai: DS
%% muc: [VD, VD, TH, VD]
%% dap_an: "SSĐS"
%% nguon: "Ảnh giáo viên gửi 24/10/2026 (M06): Kỳ thi thử đợt IV, Thi thử TN THPT 2027 (Biên soạn: Aimbet), Phần 2 Câu 2"
%% kien_thuc: "Bài 15 (phương trình tham số của đường thẳng), Bài 16 (góc giữa đường thẳng và mặt phẳng), Bài 17 (giao của đường thẳng và mặt cầu), Bài 8 (khoảng cách, tích có hướng hoặc hàm bậc hai để tìm khoảng cách nhỏ nhất); chuyển động thẳng đều"
%% ghi_chu: "Đề gốc không có hình; đáp án tự giải"
%% trang_thai: chinh-thuc
\begin{ex}
Trong không gian với hệ trục tọa độ $Oxyz$ (đơn vị trên mỗi trục tọa độ tương ứng với $1$ km) và mặt phẳng $(Oxy)$ biểu diễn mặt đất, một trạm radar được đặt tại gốc tọa độ $O$ và là tâm của vùng phủ sóng radar có dạng mặt cầu với bán kính $R=9$ km. Vào lúc $8$ giờ sáng, trạm radar phát hiện một chiếc máy bay ở vị trí $A(1;2;10)$ bay theo hướng của vectơ $\vec v=(1;2;-2)$. Chiếc máy bay bay với vận tốc là $4$ km/phút khi nằm trong vùng phủ sóng radar và bay với vận tốc là $3$ km/phút khi nằm ngoài vùng phủ sóng radar.
\choiceTF
{Chiếc máy bay nằm trong vùng phủ sóng radar trong khoảng thời gian là $40$ giây.}
{Chiếc máy bay tiếp xúc với mặt đất vào lúc $8$ giờ $5$ phút.}
{\True Giá trị sin của góc được tạo bởi đường thẳng mà chiếc máy bay bay theo và mặt đất bằng $\dfrac23$.}
{Khoảng cách từ chiếc máy bay đến trạm radar đạt giá trị nhỏ nhất bằng $2\sqrt5$ km.}
\loigiai{Vị trí máy bay $M(1+t;2+2t;10-2t)$, $t\ge0$; $|\vec v|=3$ nên quãng đường từ $A$ đến $M$ là $3t$ km.
$OM^2=81\Leftrightarrow9t^2-30t+24=0\Leftrightarrow t=\dfrac43$ hoặc $t=2$. Quãng đường trong vùng phủ sóng $3\left(2-\dfrac43\right)=2$ km, thời gian $\dfrac24$ phút $=30$ giây: a) sai.
Chạm đất khi $10-2t=0\Leftrightarrow t=5$. Thời gian: $\dfrac43+\dfrac12+\dfrac93=\dfrac{29}6$ phút $=4$ phút $50$ giây, tức lúc $8$ giờ $4$ phút $50$ giây: b) sai.
$\sin\varphi=\dfrac{|\vec v\cdot\vec k|}{|\vec v|}=\dfrac23$: c) đúng.
$OM^2=9t^2-30t+105$ nhỏ nhất tại $t=\dfrac53\in[0;5]$, $OM^2_{\min}=80$, $OM_{\min}=4\sqrt5$ km: d) sai.}
\end{ex}
